% ============================================================================
% Fichier  : partie2/P02_C04_qui_controle_si_publics.tex
% Titre    : Qui Contrôle les SI Publics au Cameroun
% Auteur   : MINKA MI NGUIDJOÏ Thierry Emmanuel
% Version  : 0.1 -- novembre 2026
% ============================================================================
% Textes lus (dossier Normes) : décret n° 2013/287 (art. 1, 2, 3, 5, 27, 28) ;
% loi n° 2010/012 (art. 4 pt 9, art. 7 al. 2, art. 13 al. 1, art. 32).
% Chambre des Comptes, lois de 2018 : textes absents du dossier Normes ;
% aucun article cité, rôle décrit en termes généraux (décision de l'auteur :
% la Chambre des Comptes s'y retrouve « par extension »).
% Fait CAPNO établi ici : rapport d'audit interne de HAMADOU transmis au
% CONSUPE (art. 5) le 15 mai 2026.
% ============================================================================

\chapter{Qui Contrôle les SI Publics au Cameroun}
\label{chap:institutions}
\chaptermark{Qui contrôle les SI publics}

\epigraph{%
    « Le mandat précède la méthode. »%
}{Règle du cours}

\niveauChapitre{Fondamental}{Chapitre~\ref{chap:audit-public}.}

\begin{boxobjectifs}
\begin{itemize}
    \item Citer les dispositions du décret n°~2013/287 qui fondent le
          contrôle des systèmes d'information par le CONSUPE.
    \item Expliquer le rôle de l'ANTIC dans l'audit de sécurité des
          systèmes d'information publics, au titre de la loi n°~2010/012.
    \item Situer la Chambre des Comptes et expliquer pourquoi elle rencontre
          les systèmes d'information par extension.
    \item Situer les autres acteurs~: inspections générales, contrôle
          financier, audit interne, CDBF.
    \item Déterminer, pour une situation donnée, qui est compétent et
          comment les interventions s'articulent.
\end{itemize}
\end{boxobjectifs}

\bigskip

% ============================================================================
\section*{Pourquoi ce chapitre}
% ============================================================================

Dans le prologue, la vérificatrice a exigé une extraction de SIDÉPENSE
sans avoir à demander la permission. Elle le pouvait parce qu'un texte
lui en donne le droit. Ce chapitre présente ces textes. Il répond à trois
questions~: \emph{qui a compétence pour contrôler un système
d'information public~? Sur quoi porte chaque compétence~? Que devient un
constat d'un organe à l'autre~?}

Un principe de méthode gouverne le chapitre~: seuls sont cités entre
guillemets les textes lus dans leur version disponible. Les autres sont
décrits dans leurs grandes lignes, et l'auditeur doit les lire avant de
les invoquer.

\separateur

% ============================================================================
\section{Le Contrôle Supérieur de l'État}
\label{sec:c04-consupe}
% ============================================================================

\subsection{Une institution supérieure de contrôle}

Le Contrôle Supérieur de l'État est organisé par le décret n°~2013/287 du
4~septembre~2013~\cite{decret2013287consupe}. Ses services relèvent de
l'autorité directe du Président de la République, à qui ils rendent
compte, et sont dirigés par un ministre délégué à la Présidence (art.~1).

\begin{boxjuris}
\textbf{\loicam{Décret n°~2013/287}, article~2, alinéa~1}~: \textit{« Les
services du Contrôle supérieur de l'Etat constituent l'institution
supérieure de contrôle des finances publiques (ISC) du Cameroun. Ils sont
chargés de l'audit externe. »}

\medskip
\textbf{Article~2, alinéa~3}~: \textit{« Dans l'exercice de leurs
attributions, les services du Contrôle supérieur de l'Etat effectuent~:
un contrôle de conformité et de régularité~; un contrôle financier~; un
contrôle de performance~; un contrôle de l'environnement~; un contrôle des
systèmes d'information. »}

\medskip
\textbf{Article~2, alinéa~4}~: \textit{« Les contrôles sont effectués par
les vérificateurs du Contrôle supérieur de l'Etat sur pièces et sur place.
A cet effet, et nonobstant les contrôles effectués par les autres organes
de contrôle, les responsables des structures auditées sont tenus de
présenter les pièces relatives à la période auditée. »}

\medskip
\textbf{Portée pour l'audit.} Trois conséquences pratiques. Le contrôle
des systèmes d'information est une compétence \emph{explicite}, écrite à
côté du contrôle de régularité~: il n'a pas à être déduit d'une autre
mission. Il s'exerce sur place, donc dans l'application, et pas seulement
sur les documents qui en sortent. Enfin, l'obligation de présenter les
pièces vaut «~nonobstant~» les autres contrôles~: une entité ne peut pas
refuser une extraction au motif que l'ANTIC ou une inspection l'a déjà
auditée.
\end{boxjuris}

\subsection{Le champ du contrôle}

L'article~2, alinéa~1, fixe un champ large. Il couvre notamment les
services publics, les établissements publics, les collectivités
territoriales décentralisées, les entreprises publiques et parapubliques,
les projets à financements extérieurs, et les organismes qui bénéficient
de concours financiers, d'avals ou de garanties de l'État. Partout où
l'une de ces entités exploite un système d'information, le contrôle des
SI s'applique.

L'organisation des services en tire les conséquences. L'article~27,
alinéa~2, crée au sein de la Division des inspections et de contrôle des
administrations publiques une \emph{section des administrations
financières, économiques, sociales et des systèmes d'informations}.
L'article~28 charge chaque section du suivi et de l'évaluation des
activités des structures de contrôle interne
(Chapitre~\ref{chap:sci}).

\subsection{Les suites : le CDBF}

\begin{boxjuris}
\textbf{\loicam{Décret n°~2013/287}, article~3, alinéa~1}~: \textit{« Les
services du Contrôle supérieur de l'Etat concourent à la sanction des
ordonnateurs et des gestionnaires des deniers publics, dans les
conditions prévues par les lois et règlements en vigueur. A ce titre, le
ministre délégué en charge du Contrôle supérieur de l'Etat préside le
conseil de discipline budgétaire et financière, dont l'organisation et le
fonctionnement sont fixés par un texte particulier. »}

\medskip
\textbf{Portée pour l'audit.} Un constat du vérificateur peut fonder une
poursuite devant le Conseil de discipline budgétaire et financière
(\sigle{CDBF}). Il doit donc être établi avec la rigueur d'une pièce de
procédure~: critère cité, preuve conservée, contradictoire respecté
(Chapitre~\ref{chap:constater}).
\end{boxjuris}

\subsection{Le lien avec l'audit interne}

\begin{boxjuris}
\textbf{\loicam{Décret n°~2013/287}, article~5}~: \textit{« (1) Les
structures d'audit interne transmettent aux services du Contrôle supérieur
de l'Etat, à la diligence des chefs de départements ministériels
concernés, un exemplaire de chacun des rapports rédigés à la suite de
toutes les interventions menées au cours de l'exercice, ainsi qu'un
exemplaire du rapport annuel d'activités. (2) Après examen des rapports
visés à l'alinéa 1er ci-dessus, les services du Contrôle supérieur de
l'Etat peuvent, le cas échéant, saisir le conseil de discipline
budgétaire et financière. »}
\end{boxjuris}

\begin{boxscenario}{Dossier CAPNO. Un rapport qui voyage}
\textbf{Constatation.} Le 15~mai~2026, M.~HAMADOU, auditeur interne de la
CAPNO, remet son rapport sur la paie~: 37 agents payés ne figurent pas au
fichier du personnel. Conformément à l'article~5, le rapport est transmis
au CONSUPE.

\medskip
\textbf{Analyse} (confiance élevée). Le rapport de l'auditeur interne
devient une pièce que le CONSUPE peut examiner, et sur laquelle il peut
saisir le CDBF. Ses constats doivent donc être reproductibles par un
vérificateur~: extraction datée, empreinte, critère cité. Un rapport qui
écrit seulement «~des anomalies ont été relevées dans la paie~» ne pourra
pas être repris.
\end{boxscenario}

\separateur

% ============================================================================
\section{L'ANTIC : l'audit de sécurité}
\label{sec:c04-antic}
% ============================================================================

L'Agence nationale des technologies de l'information et de la
communication (\sigle{ANTIC}) est le régulateur de la sécurité
électronique. La loi n°~2010/012 relative à la cybersécurité et à la
cybercriminalité définit l'audit de sécurité et en confie la régulation à
l'Agence~\cite{loi2010012}.

\begin{boxjuris}
\textbf{\loicam{Loi n°~2010/012}, article~4, point~9}~: \textit{« Audit de
sécurité~: examen méthodique des composantes et des acteurs de la
sécurité, de la politique, des mesures, des solutions, des procédures et
des moyens mis en oeuvre par une organisation, pour sécuriser son
environnement, effectuer des contrôles de conformité, des contrôles
d'évaluation de l'adéquation des moyens (organisationnels, techniques,
humains, financiers) investis au regard des risques encourus,
d'optimisation, de rationalité et de performance. »}

\medskip
\textbf{Article~7, alinéa~2} (extrait)~: l'Agence est notamment chargée
\textit{« de s'assurer de la régularité, de l'effectivité des audits de
sécurité des systèmes d'information suivant les normes en la matière, des
organismes publics et des autorités de certification »}.

\medskip
\textbf{Article~32}~: \textit{« (1) Les réseaux de communications
électroniques et les systèmes d'information sont soumis à un régime
d'audit de sécurité obligatoire et périodique de leurs systèmes de
sécurité par l'Agence. (2) L'audit de sécurité et les mesures d'impact de
gravité sont effectués chaque année ou lorsque les circonstances
l'exigent. (3) Les rapports d'audit sont confidentiels et adressés au
Ministre chargé des Télécommunications. »}

\medskip
\textbf{Portée pour l'audit.} L'ANTIC porte sur la \emph{sécurité} des
systèmes~: protection, disponibilité, intégrité. Elle ne juge pas la
régularité des opérations financières qui y transitent. Pour le
vérificateur, l'existence et la date du dernier audit de sécurité sont
une pièce à demander~; son absence est un constat. Son contenu est
confidentiel~: l'auditeur demande au moins la date, le périmètre et le
plan d'action qui en est issu.
\end{boxjuris}

\separateur

% ============================================================================
\section{La Chambre des Comptes : par extension}
\label{sec:c04-chambre}
% ============================================================================

La Chambre des Comptes de la Cour suprême juge les comptes des comptables
publics et contrôle la gestion des entités publiques. Son texte
d'organisation ne figure pas dans le dossier des normes du cours~; il
n'est donc pas cité ici, et l'auditeur doit le lire avant de l'invoquer.

Son rapport aux systèmes d'information est indirect, mais réel. Les
comptes qu'elle juge sont produits par des applications~: un compte de
gestion édité par une application de comptabilité ne vaut que ce que
valent les contrôles de cette application. La Chambre rencontre donc les
SI \textbf{par extension}~: lorsqu'elle s'interroge sur la fiabilité d'un
compte, elle s'interroge sur le système qui l'a produit.

\begin{boxinfo}
\textbf{Ce que cela change pour le juge des comptes.} Un écart entre deux
états produits par la même application, ou entre l'application et les
relevés bancaires, n'est pas seulement une irrégularité comptable~: c'est
le symptôme d'un contrôle d'application ou d'un contrôle d'interface
défaillant (Chapitre~\ref{chap:applications}). Le juge peut demander
quel contrôle aurait dû l'empêcher, et qui en avait la charge.
\end{boxinfo}

\separateur

% ============================================================================
\section{Les autres acteurs}
\label{sec:c04-autres}
% ============================================================================

\begin{tableaucours}{L{3.2cm}L{3.8cm}L{5.2cm}}{%
    \tableauHeader{Acteur} &
    \tableauHeader{Nature du contrôle} &
    \tableauHeader{Rapport au système d'information}
}
Inspections générales (ministère des Finances, ministères techniques) &
Contrôle interne de l'administration, hiérarchique &
Contrôle sur pièces et sur place, souvent déclenché par un événement \\
Contrôle financier & Contrôle a priori de la dépense & Son visa passe de
plus en plus par l'application~: il est un profil \\
Audit interne des ministères et des entités & Évaluation du contrôle
interne pour la direction & Rapports transmis au CONSUPE (art.~5) \\
CDBF & Sanction des fautes de gestion & Juge des suites, sur la base des
constats \\
Chambre des Comptes & Juridictionnel, sur les comptes & Par extension,
via les comptes produits par les SI \\
ANTIC & Régulation de la sécurité électronique & Audit de sécurité des
SI des organismes publics \\
\end{tableaucours}
\vspace{-0.8em}
\begin{center}{\small\itshape Tableau~4.1. Les acteurs du contrôle des SI
publics (rôles généraux~; textes à lire avant citation, sauf le décret
n°~2013/287 et la loi n°~2010/012)}\end{center}

\separateur

% ============================================================================
\section{Qui fait quoi, et comment s'articuler}
\label{sec:c04-articuler}
% ============================================================================

\subsection{Trois angles sur un même système}

Pour un même système, trois questions différentes sont posées par trois
institutions.

\begin{tableaucours}{L{2.6cm}L{4.8cm}L{4.8cm}}{%
    \tableauHeader{Institution} &
    \tableauHeader{Question} &
    \tableauHeader{Exemple sur SIDÉPENSE}
}
CONSUPE & Le système et son contrôle interne garantissent-ils la
régularité, la fiabilité et la performance~? & La séparation entre
liquidation et ordonnancement est-elle imposée~? \\
ANTIC & Le système est-il sécurisé selon les normes~? & Les mots de
passe, les sauvegardes, la protection du réseau sont-ils conformes~? \\
Chambre des Comptes & Les comptes produits sont-ils réguliers et
sincères~? & Le compte de gestion reflète-t-il les mandats payés~? \\
\end{tableaucours}
\vspace{-0.8em}
\begin{center}{\small\itshape Tableau~4.2. Trois angles sur un même
système}\end{center}

\subsection{Règles d'articulation}

\begin{enumerate}
    \item \textbf{S'informer des autres interventions.} Le vérificateur
          demande, dès la prise de connaissance, les rapports d'audit
          interne (art.~5), la date du dernier audit de sécurité de
          l'ANTIC et les éventuelles procédures en cours devant la Chambre
          des Comptes ou le CDBF.
    \item \textbf{S'appuyer sans déléguer.} Un rapport d'un autre organe
          oriente les tests~; il ne remplace pas la preuve que le
          vérificateur obtient lui-même.
    \item \textbf{Rester dans son angle.} Le vérificateur qui constate une
          faille de sécurité la consigne et la signale~; il ne conduit pas
          un audit de sécurité à la place de l'ANTIC.
    \item \textbf{Préserver la suite.} Dès qu'un constat peut fonder une
          sanction ou une poursuite, la preuve doit être conservée selon
          les règles du Chapitre~\ref{chap:preuve}.
\end{enumerate}

\begin{boxalert}
\textbf{Violation typique.} Reprendre dans un rapport la conclusion d'un
audit de sécurité («~le SI est sécurisé~») comme preuve que les contrôles
d'accès de l'application sont efficaces.

\smallskip
\textbf{Règle~:} un audit de sécurité et une évaluation du contrôle
interne ne répondent pas à la même question. Chacun cite l'autre, aucun
ne remplace l'autre.
\end{boxalert}

\begin{boxmetier}{Vérificateur}
Citez toujours dans l'ordre de mission la disposition qui fonde votre
contrôle~: «~article~2, alinéa~3, du décret n°~2013/287~». Face à un
refus d'extraction, l'alinéa~4 du même article suffit.
\end{boxmetier}

\begin{boxmetier}{Juge des comptes}
Quand un compte vous paraît douteux, demandez la liste des applications
qui l'ont produit et le dernier rapport du CONSUPE ou de l'audit interne
qui les a examinées. Le défaut d'un contrôle d'application explique
souvent une irrégularité comptable.
\end{boxmetier}

\begin{boxmetier}{Contrôleur financier}
Votre visa électronique est une activité de contrôle que le CONSUPE
évaluera. Gardez la liste des personnes qui détiennent votre profil, et
signalez par écrit toute attribution que vous n'avez pas demandée.
\end{boxmetier}

\separateur

\begin{boxresume}
\begin{itemize}
    \item \textbf{Une compétence explicite.} Le décret n°~2013/287 (art.~2,
          al.~3) range le contrôle des SI parmi les cinq contrôles du
          CONSUPE, institution supérieure de contrôle chargée de l'audit
          externe.
    \item \textbf{Sur pièces et sur place.} Les pièces électroniques sont
          des pièces~; leur présentation est obligatoire, nonobstant les
          autres contrôles (al.~4).
    \item \textbf{Des suites.} Le CONSUPE préside le CDBF (art.~3) et reçoit
          les rapports des audits internes (art.~5).
    \item \textbf{L'ANTIC audite la sécurité.} Loi n°~2010/012~: audit de
          sécurité obligatoire, périodique, confidentiel.
    \item \textbf{La Chambre des Comptes, par extension.} Elle rencontre les
          SI à travers les comptes qu'ils produisent.
    \item \textbf{Trois angles, aucun ne remplace l'autre.} S'informer,
          s'appuyer sans déléguer, rester dans son angle, préserver la
          suite.
\end{itemize}
\end{boxresume}

\begin{boxexo}[title={\ding{45}\quad Exercice~4.1. Qui est compétent~? \niveaubase}]
Pour chaque situation, dites quelle institution est principalement
concernée, et sur quel texte vous vous fondez quand il est cité dans ce
chapitre.
\begin{enumerate}[(a)]
    \item Le réseau d'une recette des finances n'a jamais fait l'objet
          d'un audit de sécurité.
    \item Le compte de gestion d'un poste comptable diffère des paiements
          enregistrés dans l'application.
    \item Une DAF refuse de remettre l'extraction des mandats à un
          vérificateur, au motif qu'une inspection est en cours.
    \item L'auditeur interne d'un ministère achève un rapport sur la paie.
\end{enumerate}
\end{boxexo}

\begin{boxexo}[title={\ding{45}\quad Exercice~4.2. Lire le décret \niveaubase}]
Citez les deux dispositions du décret n°~2013/287 qui permettent à un
vérificateur (a) de conduire une mission dont l'objet est un système
d'information~; (b) d'évaluer le dispositif de contrôle interne d'un
ministère.
\end{boxexo}

\begin{boxexo}[title={\ding{45}\quad Exercice~4.3. Articuler \niveaumoyen}]
Au moment où le CONSUPE ouvre sa mission au MINÉQUIR, l'ANTIC annonce un
audit de sécurité de SIDÉPENSE. Rédigez la note de la vérificatrice à son
chef de division~: ce qu'elle demande à l'ANTIC, ce qu'elle ne fera pas
elle-même, et ce qu'elle fera pour préserver la preuve.
\end{boxexo}

\begin{boxexo}[title={\ding{45}\quad Exercice~4.4. Un rapport reproductible \niveauexpert}]
Réécrivez le constat suivant du rapport de M.~HAMADOU pour qu'un
vérificateur puisse le reprendre devant le CDBF~: «~Des anomalies
importantes ont été relevées dans la paie, notamment des agents qui ne
semblent pas exister.~» Indiquez le critère, la preuve, la méthode et ce
qui reste à établir.
\end{boxexo}

\begin{boxinfo}
\textbf{Transition.} Le mandat est posé. Reste à savoir selon quelles
règles professionnelles le conduire, et à quels critères comparer ce que
l'on trouve. Le Chapitre~\ref{chap:normes} présente les normes et les
référentiels de l'auditeur public.
\end{boxinfo}
